% ======================================================================
% Chapter 2 -- Part C: comparing designs, failure walk-through,
% decision record, end matter
% ======================================================================

\section{Comparing Designs}
\label{sec:criteria}

To compare designs fairly we need fixed yardsticks. Every case in the book uses the same eight criteria, so that designs can be compared across cases and so that nobody quietly drops the criterion on which a favorite design looks worst. Table~\ref{tab:criteria} defines each criterion by saying what a score of 1 and a score of 5 look like, and how the score will eventually be measured rather than estimated.

\begin{table}[htbp]
\centering
\small
\caption{The eight evaluation criteria.}
\label{tab:criteria}
\begin{tabularx}{\textwidth}{@{}>{\raggedright\arraybackslash}p{2.4cm}>{\raggedright\arraybackslash}X>{\raggedright\arraybackslash}X>{\raggedright\arraybackslash}p{3.0cm}@{}}
\toprule
Criterion & 1 (poor) & 5 (excellent) & Measured by \\
\midrule
Quality & Frequent wrong outcomes, some unnoticed & At or above the best current human level & Accuracy on a test set \\
Cost & Running cost per transaction erodes the value & Running cost negligible next to the value & Cost per transaction \\
Speed & No faster than today & Meets the speed target with room to spare & Time from start to finish \\
Reliability & Erratic under load or bad input & Degrades gracefully; consistent from run to run & Success rate; variation across runs \\
Security and privacy & Broad permissions; untrusted input reaches agents that can write & Least privilege; trust boundaries guarded & Permission review; attack testing \\
Control and auditability & Cannot explain why an outcome happened & Every outcome traceable to inputs, rules, sources and approvers & Share of outcomes with a full record \\
Maintainability & Changes cause surprises elsewhere & Rules, prompts, tools and models change independently & Effort for a standard change \\
Human workload & Adds work or moves it to a scarcer role & People handle only real exceptions & Minutes per transaction by role \\
\bottomrule
\end{tabularx}
\end{table}

Three remarks help in using the criteria. First, at design time most scores are estimates, made from experience and quick prototypes, and each case later replaces them with measurements. Second, the criteria pull against one another. Adding a checkpoint raises quality and control but lowers speed, and that tension is the reason trade-offs exist at all. Third, some requirements are not trade-offs but \emph{constraints}. A legal or policy requirement, such as ``finance approves credits above \$2{,}000,'' rules a design out however well it scores, so constraints are checked first and only the surviving designs are scored.

\subsection{Three designs for returns}
\label{sec:tradeoff}

Figure~\ref{fig:returns-designs} shows three designs for Harbor's returns process. All three respect the two constraints (the inspector judges misuse, and finance approves large credits). They differ in how much the system decides on its own.

The \emph{single-agent design} uses one tool-using agent with access to all three servers. It reads the request, applies the warranty rules described in its instructions, and creates the authorization and label itself. This is the do-everything agent, and it is worth drawing because most teams build it first and every sponsor will ask, ``why not just one agent?'' The \emph{reviewed-workflow design} uses a router to classify the request, code to check warranty eligibility against the table, and a single model call to draft the authorization and a reply to the customer, and then a representative approves every authorization before it is created. The \emph{straight-through design} is the same as the reviewed workflow except that the representative sees only the exceptions. Code issues the authorization automatically when the item is in warranty, its value is under \$500, the router is confident, and the customer has not been flagged as a repeat returner. Everything else goes to a representative.

\begin{figure}[htbp]
\centering
\begin{tikzpicture}[x=1cm, y=1cm, sg base/.append style={text width=1.7cm, minimum height=7.5mm, font=\scriptsize, inner sep=2pt}]
% ---- single agent
\node[font=\small\bfseries, anchor=west] at (-1,0.95) {Single agent};
\node[sg trigger] (c1) at (0,0) {Customer};
\node[sg agent] (a1) at (2.6,0) {Returns agent\\{\tiny tool loop; warranty in prompt}};
\node[sg mcp, text width=2.4cm] (m1) at (5.9,0) {ERP, warehouse, shipping\nodepart{two}{\tiny\ttfamily 9 tools incl.\ create\_rma}};
\node[sg human] (f1) at (9.1,0) {Finance\\{\tiny credits $>$\$2k}};
\draw[sg flow] (c1) -- (a1); \draw[sg call] (a1) -- (m1); \draw[sg flow] (m1) -- (f1);
% ---- reviewed workflow
\begin{scope}[yshift=-2.4cm]
\node[font=\small\bfseries, anchor=west] at (-1,0.95) {Reviewed workflow};
\node[sg trigger] (c2) at (0,0) {Customer};
\node[sg agent] (r2) at (2.3,0) {Router\\{\tiny picks a path}};
\node[sg code] (w2) at (4.6,0) {Warranty\\check};
\node[sg agent] (d2) at (6.9,0) {Draft RMA\\{\tiny one call}};
\node[sg human] (h2) at (9.2,0) {Rep approves\\every RMA};
\node[sg mcp, text width=1.9cm] (m2) at (11.6,0) {ERP, shipping\nodepart{two}{\tiny\ttfamily create\_rma, label}};
\draw[sg flow] (c2) -- (r2); \draw[sg flow] (r2) -- (w2); \draw[sg flow] (w2) -- (d2);
\draw[sg flow] (d2) -- (h2); \draw[sg call] (h2) -- (m2);
\end{scope}
% ---- straight-through
\begin{scope}[yshift=-5.0cm]
\node[font=\small\bfseries, anchor=west] at (-1,0.95) {Straight-through with exceptions};
\node[sg trigger] (c3) at (0,0) {Customer};
\node[sg agent] (r3) at (2.3,0) {Router\\{\tiny picks a path}};
\node[sg code] (w3) at (4.6,0) {Warranty \&\\safe-zone rules};
\node[sg agent] (d3) at (6.9,0) {Draft RMA\\{\tiny one call}};
\node[sg human] (h3) at (9.2,-1.1) {Rep reviews\\exceptions};
\node[sg mcp, text width=1.9cm] (m3) at (11.6,0) {ERP, shipping\nodepart{two}{\tiny\ttfamily create\_rma, label}};
\draw[sg flow] (c3) -- (r3); \draw[sg flow] (r3) -- (w3); \draw[sg flow] (w3) -- (d3);
\draw[sg call] (d3) -- node[sg lab, above] {[safe zone]} (m3);
\draw[sg flow] (w3.south) |- node[sg lab, pos=0.75, below] {[else]} (h3.west);
\draw[sg call] (h3.east) -| node[sg lab, pos=0.25, below] {[approved]} (m3.south);
\end{scope}
\end{tikzpicture}
\caption{Three designs for returns, simplified to the stage that issues the return authorization (RMA). The later inspection and credit stages are the same in the second and third designs and are omitted.}
\label{fig:returns-designs}
\end{figure}

\subsection{Scoring, weighting and checking the result}

Table~\ref{tab:tradeoff} scores the three designs. Each score has a reason, and the reasons matter more than the numbers. The single agent scores 1 on control because nobody can show afterwards which warranty rule it applied, and 5 on speed and workload because it never waits for anyone. The reviewed workflow is its mirror image. The straight-through design gives up a point of quality and a point of control compared with the reviewed workflow in exchange for two points each on speed and workload.

The weights come from the business, not from the designer's taste. Harbor's sponsor cares most about speed and correct credits, the controller adds weight to auditability, and maintainability matters least because the process rarely changes. Writing the weights down before scoring keeps them honest, and showing them to the sponsor turns a technical argument into a business conversation. The weighted score of design $d$ is
\begin{equation}
  S_d = \sum_{c=1}^{8} w_c \, s_{d,c}, \qquad \sum_{c=1}^{8} w_c = 1,
  \label{eq:weighted}
\end{equation}
where $s_{d,c}$ is design $d$'s score on criterion $c$ and $w_c$ is that criterion's weight.

\begin{table}[htbp]
\centering
\small
\caption{Trade-off table for the returns designs. Scores run from 1 to 5; the weights come from the business context.}
\label{tab:tradeoff}
\begin{tabular}{@{}lcccc@{}}
\toprule
Criterion & Weight & Single agent & Reviewed workflow & Straight-through \\
\midrule
Quality & 0.20 & 2 & 5 & 4 \\
Cost & 0.10 & 3 & 4 & 4 \\
Speed & 0.20 & 5 & 2 & 4 \\
Reliability & 0.10 & 2 & 4 & 4 \\
Security and privacy & 0.10 & 2 & 4 & 4 \\
Control and auditability & 0.15 & 1 & 5 & 4 \\
Maintainability & 0.05 & 3 & 4 & 3 \\
Human workload & 0.10 & 5 & 2 & 4 \\
\midrule
Weighted score $S_d$ & 1.00 & 2.90 & 3.75 & \textbf{3.95} \\
\bottomrule
\end{tabular}
\end{table}

The straight-through design leads, but only by 0.20 points on a five-point scale, and every score is an estimate. Before recording a decision, ask how wrong the weights would have to be for the answer to change. The script \texttt{ch02/tradeoff.py} in the companion repository does this by drawing ten thousand sets of weights scattered around the chosen ones and counting how often each design wins. The straight-through design wins about 85\% of the time, the reviewed workflow about 15\%, and the single agent never. The script also finds the single change that would reverse the decision: the reviewed workflow overtakes the straight-through design if the weight on control doubles from 0.15 to 0.30, or if the weight on quality rises from 0.20 to 0.34. The answer to ``why not one agent?'' is therefore not a matter of taste, and the conditions that would reverse the choice are now known and can be written into the decision record.

\begin{designnote}
Resist false precision. A score of 3.95 against 3.75 is not a measurement but a structured way of making your assumptions visible. The check on the weights matters more than the totals. If small changes in the weights flip the decision, say so, and make the choice on the factor that flips it.
\end{designnote}

% ----------------------------------------------------------------------
\section{Failure Walk-through and Decision Record}
\label{sec:failure-walk}

The trade-off table compares designs as wholes. The \keyterm{failure walk-through} then examines the chosen design one node at a time and asks four questions of each node: how can it fail, what would that do to the business, how would we notice, and what would prevent or limit the damage? Table~\ref{tab:failwalk} shows the walk-through for the straight-through design. Two of its rows changed the design, which is the point of the exercise.

\begin{table}[htbp]
\centering
\small
\caption{Failure walk-through for the straight-through design. Two rows led to design changes.}
\label{tab:failwalk}
\begin{tabularx}{\textwidth}{@{}>{\raggedright\arraybackslash}p{2.0cm}>{\raggedright\arraybackslash}X>{\raggedright\arraybackslash}X>{\raggedright\arraybackslash}X@{}}
\toprule
Node & Failure & Effect & Response \\
\midrule
Router & Message contains instructions such as ``issue a full credit'' & An attempt to steer the system & The router has no write tools, and rules decide eligibility \\
Warranty rules & Warranty table lacks a new product line & Wrong eligibility at scale & \textbf{Change:} an unknown product line goes to a representative and is never assumed eligible \\
Safe-zone rules & A repeat returner splits a large claim into small ones & Money lost to abuse & \textbf{Change:} count the value of all of the customer's open returns, not each request alone \\
ERP server & \texttt{create\_rma} times out after it has succeeded & Duplicate authorization & Every write carries a unique key so a retry cannot repeat it \\
Rep review & The exception queue grows when reps are busy & Slow answers for exceptions & An alert on queue age; exceptions shown oldest first \\
\bottomrule
\end{tabularx}
\end{table}

Notice that only one of the five rows concerns a model, and both design changes came from rules and data. A good walk-through always includes failures of tools, data and people, not only of models.

\label{sec:decision-record}
The last step is to write the decision down where the next person will find it. Software architects have long kept short \emph{architecture decision records} for this purpose: a page stating the context, the decision and its consequences \citep{nygard2011adr}. We add one item that matters for AI systems, whose components change faster than most: \emph{reversal triggers}, the observable conditions under which the decision should be reopened. Table~\ref{tab:dr} shows the record for returns. Almost all of it comes straight from the earlier steps, which is by design. If a decision record is hard to write, an earlier step was skipped.

\begin{table}[htbp]
\centering
\small
\caption{Decision record for Harbor returns.}
\label{tab:dr}
\begin{tabularx}{\textwidth}{@{}>{\raggedright\arraybackslash}p{2.8cm}>{\raggedright\arraybackslash}X@{}}
\toprule
\multicolumn{2}{@{}l}{\textbf{Issue low-risk return authorizations automatically; send exceptions to a representative}} \\
\midrule
Status and owner & Accepted for the pilot. Owner: the returns manager. \\
Context & 1{,}100 returns a month; three days to authorization; 6\% of credits out of warranty. Constraints: the inspector judges misuse, and finance approves credits over \$2{,}000. \\
Options considered & Single agent; reviewed workflow; straight-through with exceptions (Figure~\ref{fig:returns-designs}). \\
Decision & Straight-through with exceptions, plus the two changes from the walk-through. \\
Reasons & Highest weighted score (3.95, against 3.75 and 2.90), and it wins in about 85\% of alternative weightings. The single agent was rejected because its warranty decisions cannot be audited. \\
Consequences & Representatives handle only exceptions, expected to be 30--40\% of volume. A weekly sample of automatic authorizations is reviewed. The warranty table becomes a governed configuration with an owner. \\
Reversal triggers & Reopen the decision if out-of-warranty authorizations in the safe zone exceed 1\% in any month, if an audit finding doubles the weight on auditability, or if exceptions exceed 60\% of volume, which would erase most of the gain over the reviewed workflow. \\
\bottomrule
\end{tabularx}
\end{table}

% ----------------------------------------------------------------------
\section{Putting the Toolkit Together}
\label{sec:toolkit-together}

Every case chapter uses the toolkit in the same order, and the list below doubles as a checklist for design reviews. Run the seven-step method and name the business measures and risks. Draw at least two alternative solution graphs; in Part~II the first alternative is the Google sample as it stands. Check each graph against the eight checks, and drop any design that breaks a policy or legal constraint. Write the weights from the business context, score the surviving designs with a reason for each score, and check how sensitive the result is to the weights. Walk the chosen design node by node and fold the changes back into the graph. Finally, write the decision record and a node card for every agent, MCP server and guardrail.

Chapter~\ref{ch:business} adds the business layer that this chapter took as given: where the business measures come from, how to estimate the value of moving them, and how to put the cost of errors into money, so that several of the scores in a trade-off table can be replaced by dollars.

% ----------------------------------------------------------------------
\section*{Chapter Summary}
\addcontentsline{toc}{section}{Chapter Summary}

\begin{itemize}[itemsep=2pt]
  \item A fixed notation makes the kind, owner and purpose of every component visible, makes omissions stand out, and lets designs be compared. It has nine kinds of node, five kinds of arrow, ownership and trust boundaries, and trace tags that name the business measure or risk each node serves.
  \item Eight simple checks catch the commonest structural mistakes before anyone discusses whether a design is good.
  \item The seven-step method works from the business inward: goal and measures, people and systems, judgment or rule, knowledge, actions, human checkpoints, and failure points.
  \item Retrieve-then-answer and the approval gate appear in almost every case. The do-everything agent, policy in the prompt and the silent write are the commonest anti-patterns.
  \item Designs are scored on eight fixed criteria, with weights from the business. Checking how sensitive the result is to the weights matters more than the totals.
  \item The failure walk-through examines tools, data and people as well as models, and it usually changes the design. A decision record captures the choice, the reasons and the conditions for reopening it.
\end{itemize}

\section*{Key Terms}
\addcontentsline{toc}{section}{Key Terms}

\begin{multicols}{2}
\raggedright\small
anti-pattern \\
constraint (versus criterion) \\
decision record \\
design pattern \\
evaluation criteria \\
failure walk-through \\
handoff \\
node card \\
ownership boundary \\
reversal trigger \\
solution graph \\
tool call \\
trace tag \\
trade-off table \\
trust boundary \\
well-formed graph
\end{multicols}

\section*{Review Questions}
\addcontentsline{toc}{section}{Review Questions}

\begin{enumerate}[itemsep=2pt]
  \item Why does the notation draw a company system separately from the MCP server in front of it?
  \item Explain the difference between a tool call and a handoff. Give an example in which choosing the wrong one would leave a case with no owner.
  \item Which of the eight checks does Team A's design from Chapter~\ref{ch:foundations} fail? Give one sentence for each.
  \item In the returns case, why is warranty eligibility treated as a rule even though representatives found it hard?
  \item When is an approval gate the wrong choice, and what would you use instead?
  \item What is the difference between a constraint and a criterion, and why are constraints checked first?
  \item What is a reversal trigger, and why are reversal triggers especially important for AI systems?
\end{enumerate}

\section*{Exercises}
\addcontentsline{toc}{section}{Exercises}

\paragraph{Design track}
\begin{enumerate}[itemsep=2pt]
  \item \textbf{Complete the returns graph.} Extend the straight-through design to include the inspection and credit stages, and check the full graph against the eight checks.
  \item \textbf{Re-weight.} Suppose Harbor's auditors require every return authorization to have a named human approver. Is this a constraint or a change of weights? Redo the decision and update the decision record.
  \item \textbf{Your own process.} Apply the seven-step method to the process you broke down in Chapter~\ref{ch:foundations}'s exercises, draw two alternatives, and complete a trade-off table.
\end{enumerate}

\paragraph{Build track}
\begin{enumerate}[resume, itemsep=2pt]
  \item \textbf{Sensitivity, extended.} Modify \texttt{ch02/tradeoff.py} so that it varies the \emph{scores} as well as the weights, for example by moving each score up or down by one with some probability. Does the straight-through design remain the choice?
\end{enumerate}

% ----------------------------------------------------------------------
\section*{Lab 2: Measuring Two Designs}
\addcontentsline{toc}{section}{Lab 2: Measuring Two Designs}

\paragraph{Goal.} Replace three of the estimated scores in Table~\ref{tab:tradeoff}, for quality, cost and human workload, with measurements, by running the intake stage of the reviewed-workflow and straight-through designs on the same sample of returns.

\paragraph{Time.} About 90 minutes.

\paragraph{Steps.}
\begin{enumerate}[itemsep=2pt]
  \item Open \texttt{ch02/} in the companion repository and install the pinned requirements. The folder contains 200 synthetic return requests with their correct outcomes, a warranty table, and stand-in MCP servers for the ERP and shipping systems.
  \item Run the provided router agent on all 200 requests with \texttt{python run\_router.py}, and record its accuracy and its confidence score for each request.
  \item Write the safe-zone rule from the decision record as a plain Python function in \texttt{safe\_zone.py}, including the two changes from the failure walk-through. Write unit tests for it that do not involve the model.
  \item Run \texttt{python simulate.py --design reviewed} and \texttt{--design straight}. The script assumes a representative takes 4 minutes per review and reports, for each design, the share of correct outcomes, the model cost per request, and representative minutes per 100 requests.
  \item Choose the router-confidence threshold for the safe zone by trying values from 0.5 to 0.95 and plotting correct outcomes against the share of requests handled automatically.
  \item Convert your measurements to 1--5 scores using Table~\ref{tab:criteria}, update \texttt{tradeoff.py}, and rerun it.
\end{enumerate}

\paragraph{Deliverable.} An updated trade-off table with the measured scores marked, the threshold plot, and a revised decision record. State whether the measurements changed the decision and, if so, which reversal trigger fired.

% ----------------------------------------------------------------------
\section*{Further Reading}
\addcontentsline{toc}{section}{Further Reading}

\begin{itemize}[itemsep=2pt]
  \item \citet{anthropic2024effective} describes the workflow and agent patterns that this chapter and later cases build on.
  \item \citet{nygard2011adr} is the short essay that popularized architecture decision records.
\end{itemize}
